Neste estudo, foi desenvolvida uma metodologia robusta para a classificação automática dos graus ISUP em imagens histopatológicas de próstata, abordando desafios relevantes, incluindo desbalanceamento de classes, ruído nos rótulos e a natureza ordinal da progressão tumoral. A introdução de uma função de perda híbrida, combinando regressão ordinal e \textit{focal loss}, melhorou o desempenho do modelo ao capturar explicitamente a relação hierárquica entre os graus de malignidade. A combinação dessa estratégia com a filtragem de ruído baseada em entropia mostrou-se altamente eficaz. De forma geral, a abordagem proposta elevou o índice $\kappa_{quad}$ de 0.826 para 0.857 (aproximadamente 3\%) e a acurácia de 0.592 para 0.669 (aproximadamente 7\%) em comparação com o modelo base.

Ao alinhar a penalização de erros com a relevância clínica por meio da coerência ordinal e ao enfatizar a qualidade dos dados, o método estabelece uma base sólida para a análise patológica automatizada. O sistema demonstra forte potencial como uma ferramenta computacionalmente viável e clinicamente coerente para apoiar a tomada de decisão diagnóstica no câncer de próstata.